\documentclass[conference]{IEEEtran}
% Idioma do artigo: português do Brasil.
\usepackage[T1]{fontenc}
\usepackage[utf8]{inputenc}
\usepackage[brazil]{babel}
\usepackage{amsmath,amssymb,amsfonts}
\usepackage{graphicx}
\usepackage{booktabs}
\usepackage{xcolor}
\usepackage{url}
% Citações e referências no formato ABNT (AUTOR, ano), exigido pelo enunciado.
\usepackage[alf]{abntex2cite}

\graphicspath{{figures/}}

\begin{document}

\title{Título do Artigo}

\author{\IEEEauthorblockN{Arthur Diehl Barroso}
\IEEEauthorblockA{\textit{Departamento de Ciência da Computação} \\
\textit{Universidade de Brasília}\\
Brasília, Brasil \\
% e-mail ou ORCID
}}

\maketitle

\begin{abstract}
% Resumo: escrever ao final, após os experimentos.
\end{abstract}

\begin{IEEEkeywords}
% palavras-chave
\end{IEEEkeywords}

\section{Introdução}
% Motivação, explicação do propósito e pergunta de pesquisa.

\section{Trabalhos Relacionados}
% Revisão do estado da arte (artigos da ICLR 2026).
% (A bibliografia não pode ficar vazia: manter ao menos uma citação.)
Este trabalho parte do GmNet \cite{wang2026gmnet}.

\section{Metodologia}
% Hipótese, técnicas base e proposta de inovação.
% Indicar o repositório: \url{https://github.com/Diehlgit/ProjetoFinalRNAP_UnB2026}

\section{Experimentos}
% Organização dos experimentos e análise dos resultados.

\section{Conclusões}
% O objetivo foi atingido? Como? Reforçar o resultado principal.

\bibliography{references}

\end{document}
